% =============================================================================
\section{Introducción: del conductor al dieléctrico}\label{sec:introduccion}
% =============================================================================

En el tema anterior dimos el primer paso hacia la materia con el conductor perfecto, un cuerpo con una
reserva ilimitada de cargas que se mueven sin oposición~\cite{T3}. En el otro extremo está el
\emph{aislante perfecto}, o \emph{dieléctrico}: un material que no admite cargas libres. Sus cargas
---núcleos y electrones--- están ligadas a los átomos o moléculas y no pueden abandonar su posición,
aunque sí desplazarse ligeramente respecto de ella. Ambas idealizaciones funcionan muy bien en
electrostática, y entre las dos cubren la mayor parte de los materiales de interés.

La pregunta de este tema es qué hace un campo eléctrico con un dieléctrico. La respuesta es que lo
\emph{polariza}: cada átomo o molécula se convierte en un pequeño dipolo, o alinea parcialmente el que
ya tenía. Un dieléctrico polarizado será, a escala macroscópica, una distribución continua de dipolos,
y para describirlo bastarán la teoría del dipolo y el desarrollo multipolar de los temas
anteriores~\cite{T1,T2}. La novedad es que, como en el conductor, el campo y la polarización se
determinan mutuamente: el campo polariza el medio, y los dipolos crean a su vez un campo que se suma al
aplicado.

Seguimos el orden de los apuntes manuscritos de la asignatura y nos apoyamos en los apuntes de teoría y
en sus figuras~\cite{DFA-teoria,DFA-figuras}, que hemos recalculado y redibujado; para ampliar pueden
consultarse los textos clásicos~\cite{Griffiths,Wangsness,Purcell,Feynman,Jackson}.

\subsection{El dieléctrico ideal}\label{sec:dielectrico-ideal}

\begin{definicion}[Dieléctrico]\label{def:dielectrico}
Un \emph{dieléctrico} (o aislante perfecto) es un cuerpo material, que ocupa una región $V\subset\mathbb{R}^3$
de frontera $S=\partial V$, tal que:
\begin{enumerate}
  \item no contiene cargas libres: todas sus cargas están ligadas a átomos o moléculas, alrededor de cuya
        posición de equilibrio solo pueden desplazarse distancias del orden atómico o menores;
  \item es globalmente neutro en ausencia de cargas externas depositadas en él.
\end{enumerate}
Llamaremos \emph{cargas libres} a todas las demás: las de los conductores, los iones o electrones que
podamos depositar en el material, y en general cualquier carga cuyo movimiento no esté restringido a la
escala atómica.
\end{definicion}

\begin{observacion}[Libre y ligada no es una distinción microscópica]\label{obs:libre-ligada}
Un electrón es un electrón, y crea el mismo campo esté ligado o no. La distinción es de
\emph{control}: la carga libre es la que fijamos nosotros (depositándola, conectando una batería), mientras
que la ligada es la que el material reorganiza por sí solo en respuesta al campo. Esta asimetría es la que
hará útil el vector desplazamiento de la \cref{sec:desplazamiento}.
\end{observacion}

\subsection{Notación y convenios}\label{sec:notacion}

\begin{notacion}\label{not:tema4}
Usamos la notación de los temas anteriores~\cite{T0,T1,T2,T3} y, además:
\begin{itemize}
  \item $\vr$ es el punto de campo y $\vr\,'$ el punto fuente; $\vR=\vr-\vr\,'$, $R=\abs{\vR}$ y
        $\uv R=\vR/R$. El operador $\grad'$ actúa sobre las coordenadas primadas.
  \item $\rhof$ y $\sigmaf$ son las densidades de carga libre, volúmica y superficial; $\rhoP$ y $\sigmaP$, las de
        carga ligada o \emph{de polarización}. $\qfenc$ y $\qPenc$ son las cargas libre y ligada
        encerradas por una superficie.
  \item $\vp$ es el momento dipolar de un átomo o molécula; $\vP$, el vector polarización; $\vD$, el vector
        desplazamiento; $\chi$, $\varepsilon$ y $\er$, la susceptibilidad, la permitividad y la permitividad relativa.
  \item $\nn$ es la normal unitaria exterior al dieléctrico. En una superficie que separa dos medios, $\nn$ apunta
        del medio~1 al medio~2, y los subíndices $n$ y $t$ indican las componentes normal y tangencial.
  \item \LHI{} abrevia «lineal, homogéneo e isótropo». Mientras no se diga lo contrario,
        $\phi(\infty)=0$ y $\eo\simeq\SI{8.854e-12}{F/m}$~\cite{CODATA}.
\end{itemize}
\end{notacion}
